\begin{abstract}
\noindent A smart meter that reports something surprising is telling one of two stories: either the
\emph{grid} did something unusual (an outage, a voltage sag, a demand surge) or the \emph{data} is wrong (a
drifting sensor, a corrupted value, a bypassed or falsified meter). The two need opposite responses, and
confusing them is costly. We propose a human--AI system for a smart-metered low-voltage neighbourhood that
(1) detects anomalous and biased readings with a hybrid of physics rules, statistics and unsupervised
learning, (2) explains each verdict with a calibrated confidence and an auditable record, (3) routes every
case to the right level of human control, never letting confidence override consequence, and (4) lets
several neighbourhoods learn from each other without sharing raw readings. Phase~1 delivers a labelled,
Tunisia-calibrated simulated dataset, the design of all four goals, and a first detection prototype run on
one neighbourhood: at 2.3\% false positives it detects 30 of 36 genuine events, 23 of 30 meter faults and
6 of 8 falsifications, and it shows plainly which cases remain hard.
\end{abstract}
